\section{Cuántos datos se necesitan: Scaling law y costo del ciclo de vida}

Antes de decidir las fuentes de datos, primero hay que estimar cuántos tokens se necesitan. El valor de la Scaling law es ofrecer curvas empíricas para repartir el presupuesto; su limitación es que el optimum de la loss de entrenamiento no coincide automáticamente con el optimum del costo de despliegue, de las restricciones de datos ni del domain-specific performance.

\subsection{Tres preguntas: cuántos, de dónde y cómo filtrar}

La sección anterior planteó el modelo, las GPU y los datos como un mismo problema con restricciones; esta sección empieza a desglosar la variable de los datos. Al leer la figura, las tres preguntas deben verse como un pipeline con dependencias: la cantidad de datos determina la escala de la collection, la fuente determina la provenance y el domain, y el filter determina la quality/diversity; las tres interactúan entre sí, y no se puede fijar una primero y tratar las demás como simples pasos de limpieza.

\lecturefigure{slide-13.jpg}{Las tres preguntas sobre los datos de entrenamiento: how much, where, how to filter}{deck oficial de Loubna Ben Allal, página 13}

\begin{importantbox}{Hay que distinguir raw data, retained data y training tokens}
Los raw bytes son el volumen recolectado; los retained bytes son los archivos que quedan después de filter/dedup; los unique tokens son el conteo sin duplicados después del tokenizer; los training tokens incluyen además los epochs/repeats y el mixture sampling. Si un reporte dice «3T tokens» sin aclarar a qué nivel se refiere, no se puede reproducir.
\end{importantbox}

\subsection{Scaling law: repartir el compute entre parámetros y tokens}

Determinar «cuántos datos se necesitan» no puede basarse solo en la capacidad de disco disponible; hay que colocar al mismo tiempo el tamaño del modelo, los tokens de entrenamiento y el presupuesto dentro de la curva empírica. Esta sección presenta primero la forma de la scaling-law y luego usa las conclusiones distintas de Kaplan y de Chinchilla para mostrar que el llamado optimum es una asignación de recursos obtenida bajo un conjunto de supuestos y un fitted range, no una constante independiente de los datos y del hardware.

La empirical law de estilo Kaplan suele escribirse como
\[
L(N,D)=L_\infty+aN^{-\alpha}+bD^{-\beta},
\]
donde $N$ es el número de parámetros, $D$ son los tokens de entrenamiento y $L_\infty$ es la loss irreducible. Para un Transformer dense, los FLOPs de entrenamiento suelen aproximarse como $C\approx 6ND$; con $C$ fijo, hay que elegir la combinación de modelo y datos.

\lecturefigure{slide-14.jpg}{How much data: usar scaling laws para repartir el compute budget}{deck oficial de Loubna Ben Allal, página 14}

\lecturefigure{slide-15.jpg}{Kaplan y Chinchilla: conclusiones compute-optimal distintas y ejemplos}{deck oficial de Loubna Ben Allal, página 15}

Los primeros resultados de Kaplan favorecían modelos más grandes con menos tokens; Chinchilla, con experimentos más amplios, los corrigió hacia un crecimiento más equilibrado entre parámetros y tokens. La slide compara GPT-3 175B con alrededor de 300B tokens y Chinchilla 65B con alrededor de 1.4--1.6T tokens; la conclusión didáctica es: \textbf{con el mismo compute, los modelos anteriores podían tener demasiados parámetros y estar subentrenados}.

\begin{warningbox}{Una razón constante no es una ley permanente}
«Alrededor de 20 tokens por parámetro» depende de la familia de modelos, los datos, el optimizer, el context y el fitted range. Los modelos posteriores suelen hacer overtrain de forma intencional (intentionally overtrain), y los cambios de domain o de data quality también desplazan el optimum.
\end{warningbox}

\teachervoice{00:06:20--00:09:10, la ponente usa Kaplan/Chinchilla para explicar que la scaling law es una herramienta de asignación de presupuesto, y advierte que estas conclusiones provienen de artículos y datos específicos, por lo que no conviene memorizar un único número de token-per-parameter.}

\subsection{Cómo leer las curvas de Chinchilla: cada compute budget tiene un mínimo}

La sección anterior presentó la fórmula abstracta; esta pasa a las curvas empíricas para mostrar cómo aparece el optimum en la gráfica. Al leer la figura, fije primero un color: el eje horizontal es el número de parámetros y el vertical es la validation loss; cada curva primero baja y luego sube, porque un modelo demasiado pequeño carece de capacidad y uno demasiado grande, con presupuesto fijo, no recibe suficientes tokens; el mínimo es el training-compute optimum para ese budget.

\lecturefigure{slide-16.jpg}{Chinchilla empirical curves: el mínimo en número de parámetros con compute fijo}{deck oficial de Loubna Ben Allal, página 16}

Si $C=6ND$, al sustituir $D=C/(6N)$ en la law empírica se pueden obtener las leyes de potencias de $N^*(C)$ y $D^*(C)$. Un entrenamiento real no se limita a resolver la fórmula: también debe considerar el discrete hardware shape, la parallel efficiency, el batch, el context y la disponibilidad de datos.

\begin{knowledgebox}{Cuatro cosas que la curva no dice}
Por lo general no contabiliza el serving cost; no garantiza que las downstream tasks mejoren al mismo ritmo; no describe las license/quality constraints; y tampoco incluye el costo de las fallas de entrenamiento (failure), la checkpoint selection ni la hyperparameter search.
\end{knowledgebox}

\subsection{Tres líneas de scaling evidence: Compute, Data y Model}

Las curvas de Chinchilla muestran el mínimo dentro de un presupuesto fijo; esta página va más allá y grafica la loss por separado contra el compute, el data size y el model size, mostrando los intervalos donde se ajusta una recta. Al leer la figura, confirme primero si las demás variables están cerca del óptimo y revise el rango del eje horizontal; de lo contrario, aumentar un solo eje puede limitarse a empujar el sistema a un régimen de under-training o data-limited.

\lecturefigure{slide-17.jpg}{Scaling evidence: variación de la loss como ley de potencias del compute, los data y el model size}{deck oficial de Loubna Ben Allal, página 17}

\begin{warningbox}{Las tres rectas no se pueden combinar en una causalidad incondicional}
Estas gráficas son relaciones marginales empíricas. Si la data quality baja con la escala, si la tasa de repetición aumenta o si el modelo sale del intervalo de estabilidad de la optimización, la misma pendiente no se prolonga indefinidamente. Toda predicción debe reportar el fitted range y la confidence.
\end{warningbox}

\subsection{De Chinchilla al overtraining: qué cambió en el optimum}

La línea de tiempo muestra la tendencia de Chinchilla 65B/1.4T, LLaMA 7B/1T, LLaMA-3 8B/15T, entre otros. Los equipos entrenan deliberadamente modelos más pequeños con más tokens, porque el costo adicional de un solo entrenamiento se compensa con menor inference memory, latency y deployment complexity a largo plazo.

\lecturefigure{slide-18.jpg}{Scaling-law timeline: de compute-optimal al overtraining de modelos pequeños}{deck oficial de Loubna Ben Allal, página 18}

Se define el costo del ciclo de vida del modelo como
\[
C_{\text{life}}=C_{\text{train}}+Q\cdot C_{\text{infer/token}}+C_{\text{memory}}+C_{\text{ops}},
\]
donde $Q$ es el total de tokens durante el periodo de servicio. Si $Q$ es muy grande, entrenar más un smaller model puede reducir el costo total.

\teachervoice{00:09:10--00:11:10, la ponente dice que la Chinchilla scaling law no incluye el inference cost en su objetivo; por eso hoy se eligen modelos más pequeños entrenados por más tiempo, pagando una vez un costo mayor durante el entrenamiento.}

\subsection{De GPT-1 a GPT-4: la scale amplía tanto la factura del entrenamiento como la del servicio}

La discusión anterior sobre el optimum sigue siendo abstracta; aquí se usan los cambios de orden de magnitud de la serie GPT para poner en un mismo cuadro la factura del entrenamiento y la del servicio. Las slides muestran el crecimiento de escala mediante dataset/model size y costos de entrenamiento aproximados; las cifras son estimaciones públicas, no registros contables auditados. Su función didáctica es recordar al lector que aumentar los parámetros no solo afecta a un entrenamiento, sino que encarece a largo plazo la latency por token, la GPU memory y el servicio concurrente.

\lecturefigure{slide-19.jpg}{De GPT-1 a GPT-4: crecimiento del dataset, el model size y el costo de entrenamiento}{deck oficial de Loubna Ben Allal, página 19}

\lecturefigure{slide-20.jpg}{Inference gets more expensive: número de parámetros, memoria de GPU y necesidad de GPU}{deck oficial de Loubna Ben Allal, página 20}

Si la precisión de los pesos es de $b$ bytes, la memory de una sola copia de los pesos es al menos
\[
M_{\text{weights}}=bN,
\]
y además se requieren el KV cache, las activations, el workspace y la redundancia del paralelismo. La cuantización reduce $b$, pero no reduce automáticamente todos los cuellos de botella (bottleneck) de compute y bandwidth.

\begin{knowledgebox}{Por qué se prefieren code models pequeños en el IDE}
El autocompletado de código exige baja latency, alta concurrencia y un repository context relativamente largo; cada pulsación de tecla del usuario puede generar una solicitud. Un modelo con mejor benchmark pero mayor latencia no es necesariamente el optimum del producto; esto también explica por qué StarCoder2 se ofrece en varios tamaños.
\end{knowledgebox}

\subsection{Compute-optimal no equivale a optimal: la advertencia de la Harm's law}

La sección anterior mostró que los modelos grandes pagan el inference cost una y otra vez; por eso esta sección cambia el objetivo de la loss de un solo entrenamiento al costo del ciclo de vida. Las dos slides señalan que entrenar por más tiempo un smaller model eleva el upfront compute, pero puede reducir el costo de cada deployment; la llamada ``Harm's law'' es solo un nombre humorístico para el costo adicional de entrenar más allá del punto de Chinchilla, no una nueva ley natural.

\lecturefigure{slide-21.jpg}{Compute-optimal $\neq$ deployment-optimal: por qué sobreentrenar modelos pequeños}{deck oficial de Loubna Ben Allal, página 21}

\lecturefigure{slide-22.jpg}{Harm's law: el compute adicional de entrenar más allá del punto de Chinchilla}{deck oficial de Loubna Ben Allal, página 22}

\begin{importantbox}{Para decidir, trace al menos dos curvas}
Una es training loss versus training compute; la otra es end-to-end product cost/quality versus model size. Si solo se optimiza la primera, se pasan por alto la repeated inference, la memory capacity, la energy, el edge deployment y la user latency.
\end{importantbox}

\subsection{Data constraint y quality desplazan la scaling curve}

La página de lecturas adicionales introduce el data-constrained scaling y DeepSeek LLM: cuando los unique data son limitados, repetir el entrenamiento cambia los beneficios; un corpus de mayor calidad también cambia la model/data allocation. Dicho de otro modo, $D$ no es solo un conteo de tokens: también conlleva una distribution y una quality.

\lecturefigure{slide-23.jpg}{Lecturas adicionales sobre scaling laws: data constraint y quality-dependent allocation}{deck oficial de Loubna Ben Allal, página 23}

\begin{warningbox}{El domain scaling sigue siendo un problema underexplored}
En la sesión de Q\&A, la ponente dijo que Chinchilla tiene cierto respaldo entre English/code, pero que en domains como el medical puede ser distinto; incluso dentro del mismo domain, cambiar a datos más curated puede desplazar el optimum. No copie directamente un generic web fit a un campo especializado.
\end{warningbox}

\teachervoice{00:11:10--00:13:20, la ponente presenta la data quality y el domain como caveats de la scaling-law, y señala que los resultados de DeepSeek muestran que cambiar el conjunto de datos dentro del mismo campo también modifica la asignación.}

\subsection{Resumen de la sección}

La Scaling law responde a «cómo repartir con un compute de entrenamiento fijo», pero no abarca todo el ciclo de vida del modelo. El optimum real depende también del serving volume, la memory, la latency, la unique-data constraint, el domain y la quality; por eso los equipos actuales suelen elegir modelos más pequeños con más tokens.
